% ============================================================
% OP3_1a_CV_D1_blocks_v1.tex  (EXACT INSERTION BLOCK for review;
% NOT yet implemented). Scope per GPT verdict on recon v3:
% compact "Potential-sector constraint (computed; profile
% reduction only)" paragraph next to the implemented
% OP3.1a / M-(alpha,beta) block.
% Wording precautions applied: "exact global linear-gradient
% solution of the simple P+V profile equation" (no broad
% exclusion claim); "profile-reduction descending branch" (no
% unqualified "unstable"); the xi_cond resonance is deliberately
% EXCLUDED (proposal-only per Q3); patch-extent mention is a
% hedged dimensional estimate, explicitly "not a theorem".
% No further registry echoes proposed: the existing OP-grad echo
% ("matching question recorded ... completion remains open")
% stays accurate.
% ============================================================

% ---------- ANCHOR 1 (single anchor) ----------
% Insert AFTER the end of the implemented M-(alpha,beta) paragraph:
%   "The two tensors may legitimately live at different levels; the
%    task recorded here is to exhibit the bridge explicitly.
%    OP-grad/OP3.1a sharpened; completion open."
% and BEFORE:
%   "\subsection{Collective variables and the ordered-branch infrared EFT}"

\paragraph{Potential-sector constraint (computed; profile
reduction only).}
For $S=\int d^4X\,[\,P(X)+V(\Phi)\,]$ in the profile reduction
$\Phi=\Phi(s)$, $X=\Phi'^2$, the field equation
$2\,\tfrac{d}{ds}\bigl(P_X\Phi'\bigr)=V'(\Phi)$ has the first
integral
\[
Q(X)=E+V(\Phi),\qquad Q(X):=2X\,P_X(X)-P(X),
\]
with the structural identity
$Q'(X)=P_X(X)+2X\,P_{XX}(X)$ --- exactly the longitudinal
fluctuation coefficient above; the gradient magnitude is slaved
to the potential through $dX/d\Phi=V'(\Phi)/Q'(X)$ wherever
$Q'\neq0$ and locally $\Phi'\neq0$.
Two computed consequences follow.
First, the recorded $V$ excludes an exact global linear-gradient
solution of the simple $P{+}V$ profile equation: away from
$\Phi\in\{0,\pm\phi_0\}$ one has $V'\neq0$, while the linear
profile makes the left-hand side vanish; quasi-homogeneous
gradient patches then have finite profile extent (a hedged
dimensional estimate of order $m_\sigma^{-1}$, not a theorem).
Second, the potential sector contributes only the
position-dependent mass term $\tfrac12V''(\bar\Phi)\,\chi^2$ and
the background deformation --- no derivative kinetic terms; mixed
$P(X,\Phi)$ cross-terms likewise reduce, upon integration by
parts, to background/mass-sector contributions.
Hence the bridge candidate ``coupling to the potential sector''
is reclassified: it is a constraint and a mass-sector structure,
not the kinetic bridge.
The matching question correspondingly sharpens: with
$\beta=P_X(\bar X)$ and $\beta-\alpha=Q'(\bar X)$,
\[
\alpha>\beta>0
\quad\Longleftrightarrow\quad
P_X(\bar X)>0,\qquad Q'(\bar X)<0 ,
\]
i.e.\ the broken-phase sign structure requires the operative
point to lie on a profile-reduction descending branch of $Q$,
which the simple $P{+}V$ profile mechanics does not itself
stabilise.
No inconsistency is asserted; the computation narrows where the
bridge can live (orientation/EFT-level stiffness, constraint or
composite-variable structure, or the bare-versus-EFT level
distinction).
OP-grad/OP3.1a remains open.
